\section{The knots and the volume theorem}\label{ca:intro}

For a hyperbolic knot $K$, Kashaev's volume conjecture \cite{Kashaev}
predicts
\[
 \lim_{N\to\infty}\frac{2\pi}{N}\log|K_N(K)|
      =\Vol(S^3\setminus K).
\]
We use the reduced invariant, whose value on the unknot is one.
The volume is that of the complete knot complement. The auxiliary
crossing circles introduced below are filled; the main knot cusp
remains complete. Our theorem concerns a specified class of Conway
constructions, with no bound on their fixed finite size.

\subsection{Tangles, twists, and the strictness condition}

Draw a two-strand tangle in a box with four ordered boundary points:
$0,1,2,3$ are the bottom left, bottom right, top right, and top left.
Closed components created inside a box remain part of the diagram.
An elementary augmented box contains two strands and a crossing circle
bounding a twice-punctured disk. Write $h$ for its clockwise quarter-turn.
Every occurrence is separately labeled and retains its own disk framing.

There are two operations. To form $FG$, place $F$ below $G$, join
$F_2$ to $G_1$ and $F_3$ to $G_0$, and keep the exposed tuple
$(F_0,F_1,G_2,G_3)$. The turn $\Turn$ rotates the entire box and its
disks; the new tuple is $(F_1,F_2,F_3,F_0)$.
Finally $C(FG)$ joins $0$ to $3$ and $1$ to $2$ by outside arcs.
An expression is a fixed finite recipe using these operations and
distinct elementary boxes. The order of the factors matters.

\begin{figure}[ht]
 \centering
 \resizebox{\textwidth}{!}{% Ordered boundary connections; boxes suppress their interior tangles.
\begin{tikzpicture}[x=1cm,y=1cm, every node/.style={font=\small},
 tangent/.style={line width=.7pt}, terminal/.style={circle,fill=black,inner sep=1.3pt}]
% Left: stacking.
\begin{scope}[shift={(0,0)}]
 \node[font=\small\bfseries] at (1.2,4.8) {Vertical stack $FG$};
 \draw[fill=blue!3] (.45,.45) rectangle (1.95,1.85);
 \draw[fill=blue!3] (.45,2.75) rectangle (1.95,4.15);
 \node at (1.2,1.15) {$F$}; \node at (1.2,3.45) {$G$};
 \foreach \x in {.65,1.75} {
  \draw[tangent] (\x,1.85)--(\x,2.75);
  \foreach \y in {.45,1.85,2.75,4.15} \node[terminal] at (\x,\y) {};
 }
 \node[below left] at (.65,.45) {$F_0$};\node[below right] at (1.75,.45) {$F_1$};
 \node[left] at (.65,1.85) {$F_3$};\node[right] at (1.75,1.85) {$F_2$};
 \node[left] at (.65,2.75) {$G_0$};\node[right] at (1.75,2.75) {$G_1$};
 \node[above left] at (.65,4.15) {$G_3$};\node[above right] at (1.75,4.15) {$G_2$};
\end{scope}
% Middle: a genuine clockwise turn, showing old labels in new positions.
\begin{scope}[shift={(4.25,0)}]
 \node[font=\small\bfseries] at (1.15,4.8) {Quarter-turn $\Turn F$};
 \draw[fill=blue!3] (.2,1.4) rectangle (2.1,3.3);
 \node[align=center] at (1.15,2.35) {$F$ rotated\\clockwise};
 \foreach \p in {(.4,1.4),(1.9,1.4),(1.9,3.3),(.4,3.3)} \node[terminal] at \p {};
 \node[below left] at (.4,1.4) {$F_1$};\node[below right] at (1.9,1.4) {$F_2$};
 \node[above right] at (1.9,3.3) {$F_3$};\node[above left] at (.4,3.3) {$F_0$};
 \draw[-{Latex[length=2mm]},thick] (.45,3.95) arc[start angle=130,end angle=35,radius=1.1];
 \node[align=center,font=\footnotesize] at (1.15,.55) {Labels name the old endpoints;\\the tangle rotates with them.};
\end{scope}
% Right: closure arcs are actual identity connections outside both boxes.
\begin{scope}[shift={(8.65,0)}]
 \node[font=\small\bfseries] at (1.2,4.8) {Cylindrical closure $C(FG)$};
 \draw[fill=blue!3] (.45,.65) rectangle (1.95,1.85);
 \draw[fill=blue!3] (.45,2.65) rectangle (1.95,3.85);
 \node at (1.2,1.25) {$F$};\node at (1.2,3.25) {$G$};
 \foreach \x in {.65,1.75} {
  \draw[tangent] (\x,1.85)--(\x,2.65);
  \foreach \y in {.65,1.85,2.65,3.85} \node[terminal] at (\x,\y) {};
 }
 \draw[tangent] (.65,.65)..controls(-.45,.1)and(-.45,4.4)..(.65,3.85);
 \draw[tangent] (1.75,.65)..controls(2.85,.1)and(2.85,4.4)..(1.75,3.85);
 \node[below] at (.65,.5) {$F_0$};\node[below] at (1.75,.5) {$F_1$};
 \node[above] at (.65,4.0) {$G_3$};\node[above] at (1.75,4.0) {$G_2$};
\end{scope}
\end{tikzpicture}

}
 \caption{The endpoint conventions. The boxes suppress the interior
 tangles. A turn moves the entire tangle; the labels in the middle
 panel identify its old endpoints in their new positions.}
 \label{ca:tangle-figure}
\end{figure}

The strictness test assigns a real label to the output of every
subexpression. At a leaf require $0<a<1/2$; a turn changes $x$ to
$1-x$; at a product with inputs $x,y$ and output $z$ require
\[
                \max(x,y)<z<x+y<1.
\]
At the closing root, require instead that its two input labels sum to
one: the forced last product is omitted. The expression is
\emph{strict} if these constraints have a solution. This is a finite
feasibility condition on the chosen presentation
\cite[Lemma 5.2]{Newton}; the labels are independent of both the twist
integers and the quantum level. An elementary or turned elementary
closing child is allowed.

\paragraph{An example.}
The four-piece closure $C((h_1h_2)(h_3h_4))$ is strict: give every
leaf label $3/10$ and each of the two product outputs $1/2$.
Then $3/10<1/2<3/5<1$, and the closing outputs are complementary.
The half-twist pattern $(0,1,1,1)$ gives a knot. To check this, an
even elementary piece pairs its two bottom endpoints and its two top
endpoints, while an odd piece pairs opposite corners. Stacking one
even piece with three odd pieces creates no internal closed component;
the cylindrical outside arcs join the remaining pairs into one circle.

Fix a strict expression with $b\ge2$ circles. At circle $\ell$, retain
$\epsilon_\ell\in\{0,1,-1\}$ half twists and fill the primitive slope
$\mu_\ell+p_\ell\lambda_\ell$. Here $\mu_\ell$ is the auxiliary
circle's meridian and $\lambda_\ell$ its disk longitude. Our orientation
convention makes the resulting signed crossing power
\begin{equation}\label{ca:n}
 n_\ell=\epsilon_\ell-2p_\ell.
\end{equation}
Choose the half-twist pattern so that the main strands form one
component, counting components closed inside a product as well as at
the final closure. Full twists preserve that test. Filling all auxiliary
circles therefore gives a knot, denoted $K_{\epsilon,\mathbf p}$.

At level $N\ge3$ put $d=N-1$, $q=e^{\pi\ii/N}$ and
$\kappa=q^{d^2/2}$, with $q^t=e^{\pi\ii t/N}$. Let $D$ be the oriented blackboard diagram of the filled knot, and let $w(D)$ denote its writhe. The actual reduced tensor evaluation is corrected to zero framing by
\begin{equation}\label{ca:framing}
 \zeta_{N,D}=\theta^{-w(D)},\qquad
 \theta=q^{-d^2/2}=\kappa^{-1}.
\end{equation}
For another explicitly specified ribbon framing one uses that integer framing correction instead. The single reduced cut has unknot value one. Its identification with the reduced Kashaev invariant is \cite[Theorem 10.1]{Finite}; no additional factor is assigned to an auxiliary circle.

\subsection{The asymptotic formula}

The finite evaluation of the augmented diagram, as a function of its
auxiliary colors, will be called its \emph{source}. The \emph{pure}
source has no retained half twists. Its complete-color action has a
unique strict maximum \cite[Theorem 6.4]{Newton}. Write $F_0$ for the
maximum and $a_\ell$ for its leaf coordinates. Set
\[
 X_\ell=\cot(\pi a_\ell)>0.
\]
Section~\ref{ca:strict-pure-section} constructs a holomorphic action
$\Phi_0(t)$ for small color deficits $t$, normalized by
$\Phi_0(0)=F_0$. Only $0<a_\ell<1/2$ is needed. Define its Fourier
Legendre action by
\begin{equation}\label{ca:Psi}
 \nabla\Phi_0(\xi(z))+2\pi\ii z=0,\qquad
 \Psi_0(z)=\Phi_0(\xi(z))+2\pi\ii\,\xi(z)\cdot z.
\end{equation}
All these germs are fixed by their complete values and the original logarithmic branches.

\begin{theorem}[Long twists in a fixed strict construction]\label{ca:main}
For every fixed strict construction with $b\ge2$ and retained half-twist pattern satisfying the stated knot test, there is $P_0$ such that the following holds for every fixed signed integer tuple with $|p_\ell|>P_0$ for every $\ell$. The phase
\[
 T_{\mathbf n}(z)=\Psi_0(z)+\pi\ii\sum_{\ell=1}^b n_\ell z_\ell^2
                           +2\pi\ii\sum_{\ell=1}^b z_\ell
\]
has a unique critical point $z_{\mathbf n}$ near zero. Put
$S_{\mathbf n}=T_{\mathbf n}(z_{\mathbf n})$ and
$Q_{\mathbf n}=-T_{\mathbf n}''(z_{\mathbf n})$. There is a holomorphic nonzero Fourier amplitude $B$ near zero, defined in \eqref{ca:Fourier-amplitude}, such that
\begin{align}
 K_N(K_{\epsilon,\mathbf p})
 &=\zeta_{N,D}\,\kappa^{\frac12\sum_\ell n_\ell}
 N^{3/2}D(\mathbf n)e^{NS_{\mathbf n}}
 \bigl(1+O_{\mathbf p}(N^{-1})\bigr),\label{ca:asymptotic}\\
 D(\mathbf n)&=
 \frac{(-2\ii)^b(2\pi)^{b/2}B(z_{\mathbf n})
       \prod_\ell\sin(2\pi z_{\mathbf n,\ell})}
      {\sqrt{\det Q_{\mathbf n}}}\ne0.\label{ca:D}
\end{align}
The square root is the Gaussian branch with positive real quadratic form. The formula holds through all sufficiently large integer levels, with both parities and composite levels included. In particular the invariant is eventually nonzero and
\begin{equation}\label{ca:volume}
 \lim_{N\to\infty}\frac{2\pi}{N}\log|K_N(K_{\epsilon,\mathbf p})|
 =2\pi\Re S_{\mathbf n}
 =\Vol(S^3\setminus K_{\epsilon,\mathbf p}).
\end{equation}
There is no bound on the ratios or relative signs of the filling integers. They do not grow with $N$.
\end{theorem}

\paragraph{Density among twist assignments.}
The theorem covers a set of natural density one in $\mathbb Z^b$.
Indeed, among tuples in $[-R,R]^b\cap\mathbb Z^b$, at most
\[
 b(2\lfloor P_0\rfloor+1)(2R+1)^{b-1}
\]
have some coordinate of absolute value at most $P_0$. Their proportion
tends to zero. This counts assignments for a fixed construction and
half-twist pattern. Different assignments may give the same knot; this
is not a density statement for knots ordered by crossing number. Nor
does it leave only finitely many assignments uncovered.

\subsection{Why the proof has this form}

There are three distinct calculations. First we find the asymptotic of
a positive source near its strict maximum. Next we perform the
oscillatory fillings, retaining the derivative data at repeated roots.
Finally we identify the critical action with the volume of the specified
knot complement. Positivity in the first calculation does not by itself
settle cancellation in the second.

The finite algebra permits exact preliminary reductions. A turn substitutes
$r\mapsto N-1-r$ in a wire index. Multiplication by a prescribed tangle
packet is linear in its other input, so such a corridor may be contracted
to one matrix before applying an estimate. In the sum below we first
substitute the turn relations and solve the closing constraint, whose
coefficient is $\pm1$. Section~\ref{ca:packet} records the resulting
$2b-3$ independent coordinates and their integer offsets. This exact
reduction introduces no lattice covolume.

Fourier coefficients are always taken in the normalized character
basis on the shifted root grid. Parseval's identity then bounds the
full coefficient vector, including its Hermite derivative data, in
one calculation; see Section~\ref{ca:filling}. We do not replace a
coherent contraction by the sum of its separate coefficient squares.
Nonvanishing instead comes from the positive local density and the
whole filling-sum theorem.

Here are the inputs and their precise roles. The exact tangle algebra
is \cite[Proposition 3.2, Theorem 3.3, Proposition 4.1 and
equation (4.8)]{Newton}. Its Lemma 6.1 and Theorem 6.4 give the
primitive coordinates, strict maximum, and negative Hessian.
The complementary-source bound is \cite[Theorem 1.1]{Bounds} for
$b\ge3$ and \cite[Corollary 9.1]{Bounds} for $b=2$.
We use their pure specialization. After proving both a local Fourier
expansion and a gap on the rest of the Fourier torus, we apply
\cite[Theorem 1.1]{Filling}. The marked metric and absolute
action--volume comparison are \cite[Theorem 1.3]{StrictGeometry}.

The pure source suffices even when some $n_\ell$ are odd: we expand
each full integer power $R^{n_\ell}$ in ordinary colored circles
before taking the limit. The classical shear at the end of the proof
then restores the retained half twist and the integer filling
$p_\ell$. An odd algebraic power is never assigned a half-integral
surgery slope.
